This study investigates how Large Language Models (LLMs) represent and process the ``opaque'' vigesimal (base-20) counting system of Standard French. While most modern Indo-European languages use a base-10 (decimal) system, Standard French retains vigesimal elements for numbers 70-99 (e.g., \textit{quatre-vingts} for 80, literally ``four twenties''). This structural complexity raises the question of whether LLMs map these linguistic forms to abstract numeric values consistently. To isolate this ``vigesimal tax,'' we perform a comparative analysis with the more transparent decimal systems of Swiss and Belgian French (e.g., \textit{septante}, \textit{huitante}). We evaluate three state-of-the-art models---\gpt, \claude, and \llama---on tokenization complexity, decoding, and magnitude reasoning. We find that Standard French numerals in the 70-99 range are represented by roughly 2.3 to 2.8 times more tokens than their Swiss counterparts. However, modern models have mastered decoding these numerals into digits with near 100\% accuracy, showing no direct tax on simple translation tasks. Despite this, cross-system magnitude comparison (Standard vs.\ Swiss) reveals significant reasoning gaps, with accuracy dropping to 76\% for \claude and \llama. Furthermore, we uncover a novel ``Standard Overestimation'' bias in \claude, which incorrectly claims that Standard French forms are larger than Swiss forms with higher values. These findings demonstrate that while LLMs successfully decode complex counting systems, their internal mappings to abstract number lines remain misaligned across dialects, shifting the vigesimal tax from a decoding problem to a reasoning problem in modern language models.